% !TeX root = ../main.tex

\chapter{浮动体}

\section{三线表}
\subsection{学校说明}
表中参数应标明量和单位的符号；表序及表名置于表的上方。
\subsection{模板说明}
三线表是《撰写手册》推荐使用的格式，如表~\ref{tab:exampletable}。
\begin{table}[htb]
  \centering\small
  \caption{表号和表题在表的正上方}
  \label{tab:exampletable}
  \begin{tabular}{cl}
    \toprule
    类型   & 描述                                       \\
    \midrule
    挂线表 & 挂线表也称系统表、组织表，用于表现系统结构 \\
    无线表 & 无线表一般用于设备配置单、技术参数列表等   \\
    卡线表 & 卡线表有完全表，不完全表和三线表三种       \\
    \bottomrule
  \end{tabular}
  \note{注：表注分两种，第一种是对全表的注释，用不加阿拉伯数字排在表的下边，
    前面加“注：”；第二种是和表内的某处文字或数字相呼应的注，
    在表里面用带圈的阿拉伯数字在右上角标出，然后在表下面用同样的圈码注出来}
\end{table}

编制表格应简单明了，表达一致，明晰易懂，表文呼应、内容一致。
排版时表格字号略小，或变换字体，尽量不分页，尽量不跨节。
表格太大需要转页是，需要在续表上方注明“续表”，表头页应重复排出。



\section{插图}
\subsection{学校说明}
图：精选，要具有代表性； 清楚，坐标比例要适中，同一图上不同曲线的点用不同形状标出并说明；图中的术语、符号、单位等应同文字表述所用一致；图序及图名居中置于图的正下方。
\subsection{模板说明}
有的同学可能听说“\LaTeX{} 只能使用 eps 格式的图片”，甚至把 jpg 格式转为 eps。
事实上，这种做法已经过时。
而且每次编译时都要要调用外部工具解析 eps，导致降低编译速度。
所以我们推荐矢量图直接使用 pdf 格式，位图使用 jpeg 或 png 格式。
\begin{figure}[htb]
  \centering
  \includegraphics[width=0.35\textwidth]{figures/cugbthesis-logo-center.pdf}
  \caption{图号、图题置于图的下方}
  \label{fig:logo}
  \note{注：图注的内容不宜放到图题中。}
\end{figure}

关于图片的并排，推荐使用较新的 \pkg{subcaption} 宏包，
不建议使用 \pkg{subfigure} 或 \pkg{subfig} 等宏包。
